\begin{abstract}
We evaluate the finalized Adaptive Learning Component (ALC), a structured in-model memory module that updates slot/key state during inference without changing core transformer weights. This pre-submission pass focuses on stronger evidence rather than architectural changes: multi-seed reporting (5 seeds), a stronger no-write baseline, and an additional language-shaped delayed-recall benchmark with distractors. Across 5 seeds, ALC improves core delayed-recall MSE versus a stale-memory baseline (0.632 vs 0.885) and improves threshold recall accuracy (0.113 vs 0.031), while remaining numerically stable in a 10k-update stress test (NaN rate 0) and preserving exact persistence round-trip behavior (max behavior difference 0 in the reference run). Interface training consistently reduces retrieval loss (0.092 to 0.075 mean over seeds). On the added language-shaped benchmark, ALC improves over stale-memory baseline MSE (0.951 vs 1.057) but remains behind a no-write control in this setup, indicating unresolved write-policy/interference trade-offs. We therefore position ALC as mechanism-level evidence for controllable inference-time adaptation, not as a general solution to lifelong memory or long-context modeling.
\end{abstract}
